% Generated draft: requires mathematical and source review.
\documentclass[11pt]{article}
\usepackage[margin=25mm]{geometry}
\usepackage{fontspec}
\setmainfont{Latin Modern Roman}
\usepackage{amsmath,amssymb,mathtools,unicode-math}
\setmathfont{Latin Modern Math}
\usepackage{enumitem,needspace,xurl,hyperref}
\setlength{\parindent}{0pt}
\setlength{\parskip}{5pt}
\setlength{\emergencystretch}{3em}
\newcommand{\N}{\mathbb{N}}
\newcommand{\Z}{\mathbb{Z}}
\newcommand{\Q}{\mathbb{Q}}
\newcommand{\R}{\mathbb{R}}
\newcommand{\C}{\mathbb{C}}
\newcommand{\F}{\mathbb{F}}
\newcommand{\A}{\mathbb{A}}
\newcommand{\PP}{\mathbb{P}}
\newcommand{\E}{\mathbb{E}}
\newcommand{\eps}{\varepsilon}
\newcommand{\dd}{\,\mathrm{d}}
\begin{document}
\section*{UnsolvedMath ID: 30001824}
\textbf{Source problem code:} catalogue entry
\subsection{Definitions and mathematical statement}
Let \(M\) be a smooth complex projective manifold and let \(L\) be an ample holomorphic line bundle on \(M\). Such a pair \((M,L)\) is called a polarized projective manifold. Write \(c_1(L)\in H^2(M,\R)\) for the real first Chern class, and let \(\operatorname{Aut}(M)\) denote the group of holomorphic automorphisms of \(M\). An algebraic torus is an algebraic group isomorphic to \((\C^*)^r\); maximality here means maximality under inclusion among algebraic tori in \(\operatorname{Aut}(M)\).

A Kähler form \(\omega\) represents the polarization class when \([\omega]=c_1(L)\). Its associated Kähler metric is called extremal if the \((1,0)\)-gradient of its scalar curvature \(S_\omega\) is a holomorphic vector field:
\[
\bar\partial\bigl(\operatorname{grad}^{1,0}_{\omega}S_\omega\bigr)=0.
\]

The supplied question asks whether, for every polarized projective manifold \((M,L)\) and every maximal algebraic torus \(T\subseteq\operatorname{Aut}(M)\), the following implication holds:
\[
\begin{gathered}
(M,L)\text{ is \(K\)-stable relative to }T
\\[2pt]
\Longrightarrow
\\[2pt]
\exists\,\omega\text{ a Kähler form on }M:\quad
[\omega]=c_1(L),\qquad
\bar\partial\bigl(\operatorname{grad}^{1,0}_{\omega}S_\omega\bigr)=0.
\end{gathered}
\]

\emph{Unresolved qualification.} The supplied record does not define its notion of relative \(K\)-stability. In particular, it does not specify the admissible test configurations, the modified invariant and its normalization, or the equality condition. It also places \(T\) in \(\operatorname{Aut}(M)\) without specifying preservation of the polarization or a lift of the torus action to \(L\) or a tensor power. The hypothesis above retains the source wording; these conventions require source verification before the implication is a fully specified mathematical conjecture.
\subsection{Short English statement}
Does relative \(K\)-stability of a polarized complex projective manifold, with respect to a maximal algebraic torus of automorphisms, guarantee an extremal Kähler metric in \(c_1(L)\)? The precise stability convention and the compatibility of the torus with the polarization remain unspecified in the supplied record.
\subsection{Sources}
\begin{itemize}
\item Supplied local UnsolvedMath registry record, permanent ID \(30001824\), title: \emph{Relative K-Stability and Extremal Kähler Metrics}; collection: Oberwolfach Reports Open Problems; exact source problem code: \texttt{catalogue entry}. The record gives proposed year 2011 and supplies no named proposer.
\item Supplied source citation: \emph{Differentialgeometrie im Großen (2012)}, Oberwolfach Reports, DOI 10.4171/owr/2011/33, pp. 1857--1912. Original Oberwolfach report: \url{https://doi.org/10.4171/owr/2011/33}. The citation's parenthetical year 2012 differs from the record's proposed year 2011 and the year appearing in the DOI; their relationship is not explained in the record.
\item Supplied literature source, labeled ``Mabuchi relative stability'': \url{https://doi.org/10.2969/jmsj/06620535}.
\item Supplied literature source, labeled ``Mabuchi existence results'': \url{https://arxiv.org/abs/1307.5203}.
\item Literature assessment dated exactly \texttt{2026-08-21}: the supplied curation labels the problem \texttt{partially\_solved} and states that relative \(K\)-stability has been established as a necessary condition, that strong-stability cases give existence, and that the full extremal Kähler existence conjecture in the extracted generality remains unresolved. These are supplied dated curation claims, not independently verified conclusions. The record explicitly requests independent verification of time-sensitive status and cited claims.
\item The supplied statement status is \texttt{corrected\_verified}, accompanied by a claim of manual semantic, typographic, and LaTeX review of a meaning-preserving cleanup. This is supplied editorial metadata and does not establish the mathematical assertion or resolve its stability conventions.
\end{itemize}
\end{document}
